\section{Problem formulation}
\label{sec:problem}
A target is a benchmark client with benign observations and no attack labels available to its detection pipeline. Depending on the corpus, this client is a device, host, or capture scenario. We therefore use ``target'' rather than treating every client as an independently sampled physical device. Each target has a feature space shared with its eligible peers within its study stratum. Models are not transferred between incompatible corpus feature schemas.

Let $T_N$ denote the first $N$ observations of the target's retained benign training role, $R$ its benign reference role, and $C$ its benign calibration role. A separate benign test role and held-out attack observations are used for evaluation. Peer $k$ supplies a supervised binary classifier with attack-minus-benign logit margin $m_k(x)$. The target's own supervised classifier, if constructed for its role as a peer to other targets, is excluded from its detector. The objective is to combine transferable attack evidence with target-local anomaly evidence without fitting or calibrating on the target's attack labels.

For an alert rule $A_t:\mathcal{X}_t\to\{0,1\}$, the deployment quantities of interest are
\begin{equation}
 \operatorname{TPR}_t=\Pr\{A_t(X)=1\mid Y=1,t\},\qquad
 \operatorname{FPR}_t=\Pr\{A_t(X)=1\mid Y=0,t\}.
 \label{eq:objective}
\end{equation}
The design seeks higher detection with an acceptable false-alert burden. These probabilities define the objective, not a solved constrained optimization problem: target attacks are unavailable during fitting, and the calibration parameter below does not certify a population constraint. The benchmark estimates class-specific rates on retained observations and macro-averages attack types, targets, and seeds as specified in Section~\ref{sec:experiment}.
The local fitting budget $N$ is not the total number of target observations required. Reference and calibration observations are additional. Furthermore, the executed preprocessing pools the complete benign training roles of the loaded clients to estimate feature means and scales before restricting local detector fitting to $T_N$. This includes target training observations beyond the first $N$. Consequently, the experiment measures sensitivity to the local detector fitting budget conditional on this preprocessing and on additional benign reference and calibration data. It does not establish deployment after observing only $N$ target rows.

We denote the nominal operating parameter by $\alpha$, the local reserve fraction by $r$, and the peer allocation by $B_P=(1-r)\alpha$. The local allocation parameter is $B_L=r\alpha$. These parameters specify calibration operations; they are not bounds on the false-alert probability of future traffic. In particular, the implemented resolution rule can increase the sum of the component allocations above $\alpha$, as detailed below.

The deployment motivation assumes trustworthy peer models, compatible feature extraction, and target observations correctly identified as benign. The evaluation does not test compromised peers, poisoning of benign observations, adaptive evasion, or online distribution-shift handling. Model exchange is an architectural motivation, not an experimentally established privacy guarantee: the benchmark itself prepares clients centrally and uses pooled preprocessing. No claim of federated optimization, cryptographic confidentiality, or differential privacy follows from this design.

